\chapter{Fusion spaces and the meaning of an associator}
\label{ch:fusion}

The first obstacle in learning about Haagerup symmetry is that the word
``symmetry'' no longer means a group acting by invertible matrices.  We
will replace that familiar picture one piece at a time.  This chapter
requires linear algebra and the tensor product of vector spaces.  Its
main calculation is the Fibonacci associator: a small matrix whose
entries are forced by consistency of fourfold fusion.  The calculation
will explain what an $F$-symbol contains, why its entries depend on
choices of bases, and why a fusion table is only the beginning of a
symmetry theory.

\section{Three vector spaces that should be kept separate}
\index{fusion space}

A quantum system has a Hilbert space $\mathscr H$.  A finite group $G$
acts on it through unitary operators $U_g$ satisfying
$U_gU_h=U_{gh}$.  In a two-dimensional field theory one can represent
$U_g$ by a line crossing the spatial circle.  Moving this line without
crossing an operator does not change a correlation function.  Such a
line is called \emph{topological}.  Putting two parallel lines together
realizes their product.

For a more general topological line, putting two lines together can
produce a direct sum of different line types:
\begin{equation}
 a\otimes b\simeq\bigoplus_c N_{ab}^{\ c}\,c,
 \qquad N_{ab}^{\ c}\in\Z_{\geq0}.
 \label{eq:fusion-rule-general}
\end{equation}
The labels $a,b,c$ denote simple line types.  ``Simple'' means that a
line has no nontrivial direct-sum decomposition.  The unit object
$\one$ is the invisible line.  An invertible line has an inverse under
fusion; a noninvertible line does not.

There are now three different kinds of vector spaces:
\begin{enumerate}
\item The physical Hilbert space $\mathscr H$ contains states with
energies, momenta, and possibly infinitely many descendants.
\item A fusion space
\[
 V_{ab}^{c}=\Hom(c,a\otimes b)
\]
contains junctions at which a $c$ line splits into $a$ and $b$.
Its dimension is the nonnegative integer $N_{ab}^{\ c}$.
\item A representation space of the algebra generated by closed
line operators describes how a collection of states transforms.  Its
ordinary vector-space dimension is generally unrelated to the quantum
dimension of a line.
\end{enumerate}
These distinctions remain important even when all the fusion spaces
have dimension zero or one.  Multiplicity-free fusion makes individual
junction spaces simple, but repeated fusion still creates larger spaces.

For instance, suppose a line $\tau$ obeys
\begin{equation}
 \tau\otimes\tau\simeq\one\oplus\tau.
 \label{eq:fibonacci-rule}
\end{equation}
Then $\Hom(\tau,\tau^{\otimes3})$ has dimension two.  One can fuse
the first pair to $\one$, or fuse it to $\tau$, before reaching total
charge $\tau$.  This is a two-dimensional space of junction networks;
it is not an assertion that the physical theory has two states.

\section{Semisimplicity, duality, and positivity}
\index{fusion category}

\begin{definition}
A fusion category over $\C$ is a $\C$-linear rigid monoidal category
with a simple unit, finitely many simple isomorphism classes,
finite-dimensional morphism spaces, and a decomposition of every
object as a finite direct sum of simple objects.
\end{definition}

Each word has a concrete role.  Linearity permits superpositions of
junctions. Semisimplicity means that every object decomposes into
simples; the separate finiteness assumption limits the list of
simple line types.
The monoidal product describes putting lines together.  Rigidity means
that a line has a dual $a^*$, represented physically by reversing its
orientation.  Evaluation and coevaluation junctions can create and
annihilate a line--antiline pair.  Their zigzag identities say that a
line bent into a cup and then a cap can be straightened.
\index{rigidity}\index{duality!evaluation and coevaluation}
In formulas, choose a left dual $X^*$ with maps
\[
 \operatorname{ev}_X:X^*\otimes X\to\one,
 \qquad \operatorname{coev}_X:\one\to X\otimes X^*.
\]
With associators understood, the two straightening identities are
\begin{align*}
 (1_X\otimes\operatorname{ev}_X)
 (\operatorname{coev}_X\otimes1_X)&=1_X,\\
 (\operatorname{ev}_X\otimes1_{X^*})
 (1_{X^*}\otimes\operatorname{coev}_X)&=1_{X^*}.
\end{align*}
For finite-dimensional vector spaces they say
$v=\sum_i e_i e^i(v)$, using a basis and its dual.
In a fusion category the same equations concern junction maps;
they do not require the line $X$ to be an ordinary vector space.

At the level of the fusion coefficients, rigidity implies
\begin{equation}
 N_{ab}^{\ \one}=\delta_{b,a^*}.
 \label{eq:dual-vacuum}
\end{equation}
More generally, moving a line through a cup gives Frobenius reciprocity,
for example
\[
 \dim\Hom(c,a\otimes b)=\dim\Hom(a^*\otimes c,b).
\]
The equality is an isomorphism of vector spaces obtained using actual
evaluation and coevaluation maps.  It is more than a coincidence
between integers.

A unitary fusion category also supplies adjoints of morphisms and
positive inner products.  The associativity and unit isomorphisms are
unitary, and the tensor product respects the adjoint.  In suitable
orthonormal fusion bases the recoupling matrices introduced below are
therefore unitary.  An abstract complex fusion category need not admit
this structure.

\index{quantum dimension}\index{pivotal structure}
\index{spherical structure}\index{categorical trace}
A pivotal structure is a natural monoidal identification
$j_X:X\to X^{**}$. Natural and monoidal mean that this identification
respects morphisms and tensor products. It lets us close an
endomorphism $f:X\to X$
without changing the orientation of its line. One closure gives
\[
 \Tr_{\rm cat}(f)
 =\operatorname{ev}_{X^*}
   \bigl((j_Xf)\otimes1_{X^*}\bigr)
   \operatorname{coev}_X.
\]
Sphericity says that closing on the other side gives the same scalar.
The quantum dimension is $d_X=\Tr_{\rm cat}(1_X)$.
In a unitary fusion category the canonical spherical structure has
positive dimensions $d_a$ satisfying
\begin{equation}
 d_a d_b=\sum_c N_{ab}^{\ c}d_c,
 \qquad d_{\one}=1.
 \label{eq:dimension-character}
\end{equation}
The unique positive solution is the Frobenius--Perron dimension
function. We write $\qdim(X)=d_X$ for the spherical dimension and
$\FPdim(X)$ for the positive fusion-ring dimension; they coincide
in the canonical unitary structure. A dimension is the value of a
closed line, not the number
of basis states carried by an individual microscopic site.  Irrational
values cause no conflict with the integrality of Hilbert-space
dimensions.

\begin{example}[The golden ratio from a fusion matrix]
For \eqref{eq:fibonacci-rule}, multiplication by $\tau$ on the basis
$(\one,\tau)$ has matrix
\[
 N_\tau=\begin{pmatrix}0&1\\1&1\end{pmatrix}.
\]
Its characteristic polynomial is $x^2-x-1$.  The positive eigenvalue
is
\[
 \varphi=\frac{1+\sqrt5}{2},\qquad \varphi^2=1+\varphi.
\]
The vector $(1,\varphi)^{\mathsf t}$ is its positive eigenvector.
Thus the positive dimension of $\tau$ is $\varphi$.
The other root $-\varphi^{-1}$ is also a ring character, but it cannot
be the positive dimension of a unitary line.  Later we will encounter
the Yang--Lee category, where a spherical dimension uses this other
root.  The same fusion ring supports distinct categorical structures.
\end{example}

This is an example of the Perron--Frobenius theorem. For a
nonnegative matrix whose directed graph is strongly connected, the
spectral radius $r=\max_\lambda|\lambda|$ is a simple eigenvalue
with strictly positive left and right eigenvectors, unique up to scale.
These are its \emph{Perron eigenvectors}. If some power of the
matrix has every entry strictly positive, the matrix is primitive
and every other eigenvalue has modulus less than $r$. The Fibonacci
fusion matrix above is primitive: its square has all entries positive.

\begin{example}[A trace weights the multiplicity spaces]
For $X=\bigoplus_a a^{\oplus m_a}$, an endomorphism consists of
matrices $f_a$ on the multiplicity spaces. Closing each simple line
contributes $d_a$, so
\[
 \Tr_{\rm cat}(f)=\sum_a d_a\Tr(f_a).
\]
Thus on $\one\oplus\tau^{\oplus2}$ in the Fibonacci category,
an endomorphism $(\lambda,B)$ with $B\in M_2(\C)$ has trace
$\lambda+\varphi\Tr(B)$. Its identity has trace $1+2\varphi$,
not the unweighted sum of multiplicities $3$. This weighting is the
same one that will appear in modular link traces and twist moments.
\end{example}

\section{Changing the order of fusion}
\index{associator}\index{F-symbol}

There are two natural bases of
$\Hom(d,a\otimes b\otimes c)$.  The left basis first fuses $a,b$ to
an intermediate object $e$; the right basis first fuses $b,c$ to $f$.
Suppressing multiplicity labels for the moment, write the relation as
\begin{equation}
 |((ab)e\,c);d\rangle
 =\sum_f [F^{abc}_{d}]_{ef}
 |(a\,(bc)f);d\rangle.
 \label{eq:F-basis-convention}
\end{equation}
Only admissible channels occur.  A left channel must satisfy
$N_{ab}^{\ e}N_{ec}^{\ d}\ne0$, and a right channel must satisfy
$N_{bc}^{\ f}N_{af}^{\ d}\ne0$.  Associativity of the fusion ring
guarantees equal dimensions for these two spaces.  The associator tells
us the actual change of basis.

\begin{figure}[htbp]
\centering
\begin{tikzpicture}[line width=0.8pt,scale=1.2]
\begin{scope}[xshift=-2.0cm]
 \draw (0,0)--(0,0.45)--(-0.55,1.05);
 \draw (0,0.45)--(0.85,1.7);
 \draw (-0.55,1.05)--(-1.05,1.7);
 \draw (-0.55,1.05)--(-0.05,1.7);
 \node[below] at (0,0) {$d$};
 \node[left] at (-0.3,0.78) {$e$};
 \node[above] at (-1.05,1.7) {$a$};
 \node[above] at (-0.05,1.7) {$b$};
 \node[above] at (0.85,1.7) {$c$};
\end{scope}
\begin{scope}[xshift=2.0cm]
 \draw (0,0)--(0,0.45)--(0.55,1.05);
 \draw (0,0.45)--(-0.85,1.7);
 \draw (0.55,1.05)--(0.05,1.7);
 \draw (0.55,1.05)--(1.05,1.7);
 \node[below] at (0,0) {$d$};
 \node[right] at (0.3,0.78) {$f$};
 \node[above] at (-0.85,1.7) {$a$};
 \node[above] at (0.05,1.7) {$b$};
 \node[above] at (1.05,1.7) {$c$};
\end{scope}
\draw[<->,>=Stealth] (-0.6,0.8)--node[above] {$F^{abc}_d$}(0.6,0.8);
\end{tikzpicture}
\caption{Two fusion-tree bases of the same junction space.  The
intermediate channel changes from $e$ to $f$ under recoupling.}
\label{fig:fusion-trees}
\end{figure}

When multiplicities are present, the entries of $F$ have extra indices
for bases of $V_{ab}^{e}$ and $V_{ec}^{d}$, and similarly on the right.
No new principle is involved, but one should never replace a
multidimensional junction space by a scalar merely because the line
labels are the same.

We take Hermitian inner products to be antilinear in the first
argument and linear in the second. In a unitary category,
orthonormal bases imply
\[
 F F^\dagger=F^\dagger F=I.
\]
Here and below the displayed convention treats the left channel as a
row and the right channel as a column.  A reader using the inverse
change of basis will transpose or conjugate some formulas.  Such
convention changes are harmless if carried through consistently.

\begin{example}[A local projector from an $F$-matrix]
Suppose we want to project onto right intermediate channel $f_0$,
while storing vectors in the left basis.  Taking inner products in
\eqref{eq:F-basis-convention} gives
\begin{equation}
 (P^{(f_0)})_{ee'}
 =\overline{F_{e f_0}}F_{e'f_0}.
 \label{eq:recoupling-projector}
\end{equation}
It is an outer product, hence Hermitian and positive.  Its square is
itself because the corresponding column of $F$ has norm one.  Different
channels give orthogonal projectors, and their sum is the identity.
This small formula is the basic local ingredient of an anyon-chain
Hamiltonian.
\end{example}

\section{Gauge is a choice of junction bases}
\index{gauge!fusion junctions}

In a multiplicity-free category choose a nonzero basis vector
$t_{ab}^{c}$ in every nonzero fusion space.  Replacing it by
$t_{ab}^{c\,\prime}=u_{ab}^{c}t_{ab}^{c}$ changes the two fusion-tree
bases by products of two factors.  With convention
\eqref{eq:F-basis-convention},
\begin{equation}
 [F^{abc}_{d}]'_{ef}
 =\frac{u_{ab}^{e}u_{ec}^{d}}
 {u_{bc}^{f}u_{af}^{d}}
 [F^{abc}_{d}]_{ef}.
 \label{eq:F-gauge}
\end{equation}
For orthonormal bases the $u$'s have absolute value one.  In purely
algebraic classifications any nonzero complex basis change is allowed.
For higher multiplicities the scalars become invertible matrices.

An individual entry of $F$ is therefore usually not an invariant.
Statements comparing entries must specify a gauge or exhibit a
cancellation of all gauge factors.  A useful exception occurs when
$a=b=c=d=e=f=\rho$.  The numerator and denominator of
\eqref{eq:F-gauge} are then identical.  Consequently
\begin{equation}
 [F^{\rho\rho\rho}_{\rho}]_{\rho\rho}
 \label{eq:rho-gauge-invariant}
\end{equation}
is unchanged by trivalent-basis rescaling.  We will use this elementary
invariant to distinguish two Haagerup categories with exactly the same
fusion table.

Gauge equivalence is also different from relabelling simple objects.
A tensor equivalence may include a permutation preserving the fusion
rules.  When checking whether two categories are equivalent, one must
allow both operations.  Conversely, a convenient common labelling
does not identify their associators.

\section{The pentagon is path independence}
\index{pentagon identity}

For four objects there are five binary parenthesizations.  Moving
between adjacent ones uses an associator.  Two routes from
$((a\otimes b)\otimes c)\otimes d$ to
$a\otimes(b\otimes(c\otimes d))$ must agree.  In terms of the
associator $\mathfrak a_{a,b,c}:(a\otimes b)\otimes c\to
a\otimes(b\otimes c)$, this is
\begin{align}
 \mathfrak a_{a,b,c\otimes d}\,
 \mathfrak a_{a\otimes b,c,d}
 ={}&(1_a\otimes\mathfrak a_{b,c,d})
 \mathfrak a_{a,b\otimes c,d}
 (\mathfrak a_{a,b,c}\otimes1_d).
 \label{eq:pentagon-abstract}
\end{align}
Composition is read from right to left.  In chosen fusion bases this
becomes polynomial equations among the $F$-entries, with sums over
intermediate channels.  Those equations are the pentagon equations.

\begin{figure}[htbp]
\centering
\begin{tikzpicture}[>=Stealth,semithick,every node/.style={font=\small}]
 \node (top) at (0,2.2) {$((ab)c)d$};
 \node (left) at (-3.25,0.65) {$(a(bc))d$};
 \node (bottomleft) at (-2.15,-1.25) {$a((bc)d)$};
 \node (bottomright) at (2.15,-1.25) {$a(b(cd))$};
 \node (right) at (3.25,0.65) {$(ab)(cd)$};
 \draw[->] (top)--node[above left] {$\mathfrak a\otimes1$}(left);
 \draw[->] (left)--node[left] {$\mathfrak a$}(bottomleft);
 \draw[->] (bottomleft)--node[below] {$1\otimes\mathfrak a$}(bottomright);
 \draw[->] (top)--node[above right] {$\mathfrak a$}(right);
 \draw[->] (right)--node[right] {$\mathfrak a$}(bottomright);
\end{tikzpicture}
\caption{The pentagon compares a route of three reassociations with
a route of two.  Each edge is a linear map on the relevant fusion
space; the two composite maps must agree.}
\label{fig:associator-pentagon}
\end{figure}

Why are four objects enough?  The coherence theorem for monoidal
categories says that the pentagon, together with the unit axioms,
guarantees that all reassociations of an arbitrary tensor product
agree.  One does not separately impose a new axiom for five, six, or
one hundred objects.  This is an early example of a theme that recurs
in the CFT construction: a finite consistency condition can control
arbitrarily complicated compositions, but only through a theorem
explaining the passage.

Unitarity and the pentagon are independent demands.  A random unitary
matrix does not define a fusion category.  Conversely, a solution of
the pentagon over $\C$ does not automatically supply positive inner
products making its recoupling matrices unitary.

\section{A complete small pentagon calculation: Fibonacci}
\index{Fibonacci category}

We now see the golden ratio appear a second time, from associativity
rather than the dimension character.  Use the two simple objects
$\one,\tau$ with \eqref{eq:fibonacci-rule}.  Choose strict unit
conventions and a real gauge in which the only nontrivial matrix is
\begin{equation}
 F^{\tau\tau\tau}_{\tau}
 =\begin{pmatrix}u&v\\v&-u\end{pmatrix},
 \qquad u,v>0,
 \label{eq:fib-F-unknown}
\end{equation}
in channels $(\one,\tau)$.  Take the scalar
$F^{\tau\tau\tau}_{\one}=1$.  The positive solution we obtain
will be the standard unitary Fibonacci category.

The four-$\tau$ fusion space with total charge $\tau$ has dimension
three.  In the completely left-associated tree label its channels by
\[
 (e,f)=(\one,\tau),(\tau,\one),(\tau,\tau),
\]
where $e$ is the charge after two lines and $f$ after three.  To make
every matrix reconstructible, name the four external lines $a,b,c,d$,
all equal to $\tau$.  The five trees use the following intermediate
labels:
\[
\begin{array}{c|c}
\text{tree, with total charge }\tau&\text{ordered channel pair}\\\hline
(((ab)_e c)_f d)_\tau&(e,f)\\
((a(bc)_g)_f d)_\tau&(g,f)\\
(a((bc)_g d)_h)_\tau&(g,h)\\
(a(b(cd)_k)_h)_\tau&(k,h)\\
((ab)_e(cd)_k)_\tau&(e,k)
\end{array}
\]
In every row, order the allowed pairs as
$(\one,\tau),(\tau,\one),(\tau,\tau)$.
The three-step route passes through the first four rows, and the
two-step route passes from the first to the fifth and then the fourth.
Along the three-step side the successive change-of-basis matrices are
\[
 A=\begin{pmatrix}u&0&v\\0&1&0\\v&0&-u\end{pmatrix},
 \qquad
 B=\begin{pmatrix}1&0&0\\0&u&v\\0&v&-u\end{pmatrix},
 \qquad A.
\]
Each block is simply \eqref{eq:fib-F-unknown} acting on the two
channels it can mix; the isolated entry is a one-dimensional
recoupling.  Along the two-step side the matrices are $B$ followed by
\[
 C=\begin{pmatrix}0&u&v\\1&0&0\\0&v&-u\end{pmatrix}.
\]
The permutation visible in $C$ comes from using the same final channel
order as on the other route.  Thus the pentagon is the ordinary matrix
identity
\begin{equation}
 ABA=BC.
 \label{eq:fib-pentagon-small}
\end{equation}
Multiplication gives
\[
 ABA=\begin{pmatrix}
 u^2-uv^2&v^2&uv+u^2v\\
 v^2&u&-uv\\
 uv+u^2v&-uv&v^2-u^3
 \end{pmatrix},\qquad
 BC=\begin{pmatrix}
 0&u&v\\u&v^2&-uv\\v&-uv&u^2
 \end{pmatrix}.
\]
The $(2,1)$ entry requires $v^2=u$.  Since $v\ne0$, the $(1,3)$
entry requires $u+u^2=1$.  These two relations imply every other
entry, including $u-u^3=u^2$ in the bottom right.  Positivity selects
\[
 u=\varphi^{-1},\qquad v=\varphi^{-1/2}.
\]
Therefore
\begin{equation}
 F^{\tau\tau\tau}_{\tau}
 =\begin{pmatrix}
 \varphi^{-1}&\varphi^{-1/2}\\
 \varphi^{-1/2}&-\varphi^{-1}
 \end{pmatrix}.
 \label{eq:fib-F}
\end{equation}
Its square is the identity because
$\varphi^{-2}+\varphi^{-1}=1$.  The total-charge-$\one$
four-$\tau$ pentagon reduces to this same identity.  Pentagons
containing the unit reduce to the strict unit conventions.  This
accounts for all cases in the two-object fusion rule.

This example illustrates the right order of reasoning.  We specified
the fusion spaces and a gauge ansatz, imposed associativity, selected
the positive solution, and verified unitarity.  Merely writing
$\tau^2=\one+\tau$ would not have supplied the matrix
\eqref{eq:fib-F}.

\begin{example}[The Fibonacci projector]
Projecting two neighboring $\tau$ lines onto the vacuum channel,
while keeping the left fusion basis of three lines with total $\tau$,
gives
\[
 P^{(\one)}=
 \begin{pmatrix}
 \varphi^{-2}&\varphi^{-3/2}\\
 \varphi^{-3/2}&\varphi^{-1}
 \end{pmatrix}.
\]
Its trace is one and its determinant is zero.  Its image is the span
of $(\varphi^{-1},\varphi^{-1/2})^{\mathsf t}$.  The complementary
projector is $P^{(\tau)}=I-P^{(\one)}$.  An interaction
$-P^{(\one)}$ energetically favors the vacuum fusion channel, but
overlapping interactions need not commute because the two adjacent
pairs are diagonal in different fusion bases.
\end{example}

\section{What closed line operators remember}

If $Y_a$ is the operator obtained by wrapping a simple line around
the spatial circle, then
\begin{equation}
 Y_aY_b=\sum_cN_{ab}^{\ c}Y_c,
 \qquad Y_a^\dagger=Y_{a^*}.
 \label{eq:fusion-operator-algebra}
\end{equation}
This representation remembers the fusion ring, but it does not
remember all junctions or associators.  In particular, knowing these
matrices on an untwisted Hilbert space is generally insufficient to
determine the transformation law of an operator product.  Defect
sectors and junctions enlarge the algebra to the tube algebra, studied
in Chapter~\ref{ch:centers}.

The vacuum transforms by the dimension character:
\[
 Y_a|0\rangle=d_a|0\rangle.
\]
For an invertible line $d_a=1$, so this agrees with the usual
definition of an invariant state.  A noninvertible invariant state
instead has eigenvalue $d_a$.  Requiring every line to have eigenvalue
one would contradict fusion whenever a noninvertible object has more
than one summand in its square.

\section{Exercises}

\begin{exercise}\label{ex:fib-count}
Let $m_L=\dim\Hom(\one,\tau^{\otimes L})$ and
$n_L=\dim\Hom(\tau,\tau^{\otimes L})$.  Starting with $L=0$,
derive their recursion, list the first six pairs, and explain why
their exponential growth rate is $\varphi$ although every entry is an
integer.

\par\smallskip\noindent{\footnotesize\normalfont
\hyperref[sol:ex:fib-count]{Solution on page~\pageref*{sol:ex:fib-count}.}}\par
\end{exercise}

\begin{exercise}\label{ex:gauge-cancel}
Derive \eqref{eq:F-gauge} by rescaling the two trivalent vertices
in each fusion tree.  Verify explicitly that
\eqref{eq:rho-gauge-invariant} is unchanged.  Explain why a comparison
of two other individual entries need not distinguish categories.

\par\smallskip\noindent{\footnotesize\normalfont
\hyperref[sol:ex:gauge-cancel]{Solution on page~\pageref*{sol:ex:gauge-cancel}.}}\par
\end{exercise}

\begin{exercise}\label{ex:projectors}
For an arbitrary unitary $F$-matrix, use
\eqref{eq:recoupling-projector} to prove
$P^{(f)}P^{(g)}=\delta_{fg}P^{(f)}$ and $\sum_fP^{(f)}=I$.
Keep the complex conjugates in the derivation.

\par\smallskip\noindent{\footnotesize\normalfont
\hyperref[sol:ex:projectors]{Solution on page~\pageref*{sol:ex:projectors}.}}\par
\end{exercise}

\begin{exercise}\label{ex:fib-pentagon}
Check every entry of \eqref{eq:fib-pentagon-small} from
$v^2=u$ and $u^2+u=1$.  Which step uses the assumption $v\ne0$?
Why does the negative root of $u^2+u=1$ fail the displayed positive,
real gauge?

\par\smallskip\noindent{\footnotesize\normalfont
\hyperref[sol:ex:fib-pentagon]{Solution on page~\pageref*{sol:ex:fib-pentagon}.}}\par
\end{exercise}

\begin{exercise}\label{ex:dimension-not-state}
In the Fibonacci fusion ring solve the scalar equation
$y^2=1+y$.  Interpret its two solutions as one-dimensional
representations of the complexified fusion algebra.  Explain why
neither solution is an ordinary dimension of a vector space.

\par\smallskip\noindent{\footnotesize\normalfont
\hyperref[sol:ex:dimension-not-state]{Solution on page~\pageref*{sol:ex:dimension-not-state}.}}\par
\end{exercise}

\paragraph{Source note.}
The category and coherence language is standard.  The Fibonacci
calculation is included to make the recoupling convention explicit;
the Haagerup specialization in Chapter~\ref{ch:haagerup} uses the transparent
$F$-symbol presentation of Huang and Lin.  No large table of
coefficients is needed to follow the examples in this chapter.
